%% =====================================================================
%%  structures/informe-sr/portada.tex  --  Portada, apertura de articulo,
%%  membrete en linea, separador (edicion interna) y contraportada
%% ---------------------------------------------------------------------
%%  Lo que va A SANGRE se dibuja como overlay TikZ anclado a la pagina
%%  fisica (remember picture => 2 pasadas, las de build.sh). Coordenadas
%%  en cm desde la esquina inferior izquierda del A4 (21.0 x 29.7).
%%
%%  Portada v2 (pedido de Edison 2026-09-03): SOBRIA. Fondo blanco, una
%%  sola fotografia a sangre horizontal (3:2, 21 x 14 cm), titulo grande,
%%  numero y periodo, ambito, identificacion del despacho y emblemas.
%%  Sin banda de color, sin franja textil, sin cifras, sin consigna.
%%
%%     \portadainforme[foto=, kicker=, titulo=, ambito=, codigo=]
%%     \articulo{kicker}{Titular}{Entrada}     apertura de cronica en flujo
%%     \membreteinforme                        membrete institucional en linea
%%     \capitulo[foto]{Titulo}{Subtitulo}      separador a pagina completa (edicion interna)
%%     \contraportadainforme
%% =====================================================================

%% --- Portada -----------------------------------------------------------
\pgfkeys{/isr/portada/.is family, /isr/portada,
  foto/.store in=\isrPortadaFoto,
  kicker/.store in=\isrPortadaKicker,
  titulo/.store in=\isrPortadaTitulo,
  ambito/.store in=\isrPortadaAmbito,
  codigo/.store in=\isrPortadaCodigo,
  default/.style={foto=, kicker=Informe de Semana de Representación,
    titulo={Semana de\\Representación}, ambito=, codigo=}}

\newcommand{\portadainforme}[1][]{%
  % OJO pgfkeys recorta UN solo espacio inicial por clave: sin espacio antes de #1.
  \pgfkeys{/isr/portada,default,#1}%
  \thispagestyle{empty}%
  \begin{tikzpicture}[remember picture, overlay]
    \begin{scope}[shift={(current page.south west)}]
      \fill[white] (0,0) rectangle (\paperwidth,\paperheight);
      % --- Linea institucional superior ------------------------------------
      \node[anchor=west, inner sep=0pt] at (2.0,28.30) {%
        \fuentedatossemi\fontsize{6.8}{8}\selectfont\color{colorisrtextosec}%
        \addfontfeatures{LetterSpace=16}CONGRESO DE LA REPÚBLICA · CÁMARA DE DIPUTADOS};
      \node[anchor=east, inner sep=0pt] at (19.0,28.30) {%
        \fuentedatossemi\fontsize{6.8}{8}\selectfont\color{colorisrtextosec}%
        \addfontfeatures{LetterSpace=16}PERIODO PARLAMENTARIO \meta{periodo}};
      % --- Kicker, titulo, ambito -----------------------------------------
      \node[anchor=north west, inner sep=0pt, text width=17cm, align=left] at (2.0,26.05) {%
        {\fuentedisplaysemi\fontsize{9.5}{12}\selectfont\color{colorisrprimario}%
          \addfontfeatures{LetterSpace=14}\MakeUppercase{\isrPortadaKicker}}\\[17pt]
        {\fuentedisplaynegra\fontsize{46}{49}\selectfont\color{colorisrtexto}\isrPortadaTitulo}\\[15pt]
        {\fontsize{13}{16}\selectfont\color{colorisrtextosec}\isrPortadaAmbito}};
      % --- Fotografia principal a sangre horizontal (21 x 14 cm) ------------
      \begin{scope}
        \clip (0,5.30) rectangle (\paperwidth,19.30);
        \if\relax\detokenize\expandafter{\isrPortadaFoto}\relax
          \fill[neutral200] (0,5.30) rectangle (\paperwidth,19.30);
          \node at (10.5,12.30) {%
            \fuentedatossemi\fontsize{7.5}{10}\selectfont\color{colorisrtextosec}%
            \addfontfeatures{LetterSpace=14}FOTOGRAFÍA PRINCIPAL
            \if\relax\detokenize\expandafter{\isrPortadaCodigo}\relax\else
              \ \textperiodcentered\ \isrPortadaCodigo\fi};
        \else
          \node[inner sep=0pt] at (10.5,12.30) {\includegraphics[width=\paperwidth]{\isrPortadaFoto}};
        \fi
      \end{scope}
      % --- Identificacion del despacho y emblemas ---------------------------
      \node[anchor=west, inner sep=0pt] at (2.0,3.10) {%
        \fuentedisplaynegra\fontsize{11.5}{14}\selectfont\color{colorisrtexto}%
        \addfontfeatures{LetterSpace=3}\MakeUppercase{\meta{autor}}};
      \node[anchor=west, inner sep=0pt] at (2.0,2.40) {%
        \fontsize{9.5}{12}\selectfont\color{colorisrtextosec}%
        \meta{cargo}\ifmetadef{distrito}{ por \meta{distrito}}{}};
      \node[anchor=east, inner sep=0pt] at (19.0,2.75) {%
        \includegraphics[height=1.35cm]{escudo-peru.png}\hspace{0.34cm}%
        \includegraphics[height=1.42cm]{logo-camara-mono.png}};
    \end{scope}
  \end{tikzpicture}%
  \null\clearpage
}

%% --- Apertura de cronica (articulo de revista, en flujo) ----------------
%%  \articulo{kicker}{Titular}{Entrada}: entra al indice y a los
%%  marcadores, fija la marca de cabecera. La foto de apertura, si la hay,
%%  se coloca ANTES con \isrfotoapertura (componentes.tex).
\newcommand{\articulo}[3]{%
  \par\Needspace{16\baselineskip}\vspace{14pt}%   % \Needspace (no \needspace): sin penalizacion blanda, que a una columna separaba foto y titular
  \phantomsection\addcontentsline{toc}{section}{#2}%
  \markboth{\MakeUppercase{#1}}{}%
  \noindent{\isrtxkicker\color{colorisrprimario}\MakeUppercase{#1}\par}%
  \vspace{7pt}%
  \noindent{\fuentedisplaynegra\fontsize{21}{24}\selectfont\color{colorisrtexto}\isrsinguiones#2\par}%
  \vspace{7pt}%
  \noindent{\fontsize{11.5}{15}\selectfont\itshape\color{colorisrtextosec}\isrsinguiones#3\par}%
  \vspace{10pt}}

%% --- Membrete institucional EN LINEA -----------------------------------
%%  Misma anatomia y colores que membretes/despacho (logos | caja negra
%%  del nombre | lema | filete negro), compuesto en el cuerpo para no
%%  ensanchar la cabecera de todo el documento. Solo en la presentacion.
\newcommand{\membreteinforme}{%
  \noindent\begin{minipage}{\linewidth}
    \begin{minipage}[c]{0.44\linewidth}\raggedright
      \includegraphics[height=1.50cm]{logo-camara-diputados.png}%
      \hspace{0.30cm}%
      \includegraphics[height=1.42cm]{escudo-peru.png}%
    \end{minipage}%
    \hfill
    \begin{minipage}[c]{0.54\linewidth}\raggedleft
      \setlength{\fboxsep}{0pt}%
      \colorbox{colormembrete}{%
        \begin{minipage}[c][1.30cm][c]{\despachoAnchoCaja}%
          \centering\color{white}%
          {\fuentemembretesemi\fontsize{10.5}{13}\selectfont\MakeUppercase{\meta{autor}}}\\[2.5pt]
          {\fuentemembrete\fontsize{9}{11}\selectfont\MakeUppercase{\meta{cargo}}}%
        \end{minipage}}\\[4.5pt]
      \makebox[\despachoAnchoCaja][c]{%
        {\itshape\fontsize{7.5}{9}\selectfont\color{neutral700}``\meta{lema}''}}%
    \end{minipage}\par\vspace{5pt}
    {\color{colorfilete}\rule{\linewidth}{0.6pt}}\par
  \end{minipage}\par\vspace{16pt}}

%% --- Separador de capitulo (pagina completa en vino) -- EDICION INTERNA --
\newcommand{\capitulo}[3][]{%
  \clearpage\thispagestyle{empty}%
  \phantomsection\addcontentsline{toc}{section}{#2}%
  \markboth{\MakeUppercase{#2}}{}%
  \begin{tikzpicture}[remember picture, overlay]
    \begin{scope}[shift={(current page.south west)}]
      \fill[colorisrprofundo] (0,0) rectangle (\paperwidth,\paperheight);
      \node[anchor=north west, inner sep=0pt, text width=15.5cm, align=left] at (2.0,22.0) {%
        {\fuentedisplaynegra\fontsize{28}{32}\selectfont\color{white}#2\par}
        \vspace{9pt}
        {\fontsize{12}{15.5}\selectfont\color{puka100}#3\par}};
      \if\relax\detokenize{#1}\relax\else
        \begin{scope}
          \clip (0,0) rectangle (\paperwidth,11.6);
          \node[inner sep=0pt] at (10.5,5.8) {\includegraphics[width=\paperwidth]{#1}};
        \end{scope}
      \fi
    \end{scope}
  \end{tikzpicture}%
  \null\clearpage
}

%% --- Contraportada -----------------------------------------------------
%%  Fondo vino (color profundo de la marca), retrato OFICIAL de cuerpo
%%  (\despachoRetratoCuerpo, lo fija el perfil del despacho), bloque de
%%  identificacion y contacto, emblemas. Sin franja textil ni consigna.
\newcommand{\contraportadainforme}{%
  \clearpage\thispagestyle{empty}%
  \begin{tikzpicture}[remember picture, overlay]
    \begin{scope}[shift={(current page.south west)}]
      \fill[colorisrprofundo] (0,0) rectangle (\paperwidth,\paperheight);
      \node[anchor=south west, inner sep=0pt] at (1.6,0) {%
        \includegraphics[height=13.2cm]{\despachoRetratoCuerpo}};
      %%  Nombre CORTO de marca (autorbanda) para que no se parta; el
      %%  nombre completo ya va en la portada y en el pie de cada pagina.
      \node[anchor=north west, inner sep=0pt, text width=9.0cm, align=left] at (9.95,17.2) {%
        {\fuentedatossemi\fontsize{7}{9}\selectfont\color{puka200}%
          \addfontfeatures{LetterSpace=16}CONGRESO DE LA REPÚBLICA\\CÁMARA DE DIPUTADOS\par}
        \vspace{10pt}
        {\fuentedisplaynegra\fontsize{22}{25}\selectfont\color{white}\isrsinguiones%
          \MakeUppercase{\ifmetadef{autorbanda}{\meta{autorbanda}}{\meta{autor}}}\par}
        \vspace{7pt}
        {\fuentedisplaysemi\fontsize{8.6}{12}\selectfont\color{puka200}%
          \addfontfeatures{LetterSpace=8}\MakeUppercase{\meta{cargo}\ifmetadef{distrito}{ POR \meta{distrito}}{}}\\
          PERIODO \meta{periodo}\par}
        \vspace{16pt}
        {\fontsize{9.5}{13}\selectfont\color{puka100}\meta{direccion}\par
          \vspace{3pt}\meta{correo}\ \textperiodcentered\ Cel.~\meta{celular}\par}
        \vspace{16pt}
        {\fontsize{8.8}{12}\selectfont\itshape\color{puka200}%
          \ifmetadef{notacontraportada}{\meta{notacontraportada}}{%
            Informe de la Semana de Representación elaborado por el
            despacho parlamentario y presentado conforme al Reglamento
            de la Cámara de Diputados.}\par}};
      \node[anchor=south east, inner sep=0pt] at (19.4,1.8) {%
        \includegraphics[height=1.40cm]{escudo-peru.png}\hspace{0.34cm}%
        \includegraphics[height=1.50cm]{logo-camara-blanco.png}};
    \end{scope}
  \end{tikzpicture}\null}
